\section{Conclusion}\label{sec:conclusion}
Which stored bits are worth protecting depends on how many of them fail. The expected loss of a faulty model splits into a term per bit, which grows like the bit-error probability and favours protecting the flips that raise the loss on average, and a term per pair of bits, which grows like its square and favours leaving a balanced remainder; a fixed ranking is optimal only while the first dominates. \Branch{R1}{The rate at which the preferred ranking switches can be computed from calibration data, and on nine ResNet-8 checkpoints it agreed with the measurement within a factor of 1.5, about 12\% late on average.}{\TBD{Stage 1 outcome in one sentence}.} Because protection mechanisms only lower the probability of an error, the same expression takes the residual error curve of each one, \Branch{R2}{and an assignment that minimizes it for the given rate was never more than half a point behind the best tested mask on ResNet-8 \TBD{}.}{\TBD{Stage 2 outcome in one sentence}.} Across six frozen INT8 models, significance-first protection tolerates a fraction of the fault rate that gradient-based masks tolerate, and at dense fault rates triplicating well-chosen bits outperforms blanket codes and checksum correction of equal storage. All results assume logical faults with idealized voters and decoders.
